% Replacement of every D=21 cusp-space/Hecke-matrix input.
% The finite arithmetic is checked in d21_local_free.py.

\begin{proposition}
Let $F=\mathbb Q(\sqrt{21})$. A product $E_2(\phi,\psi)h$, with $h$
a full-level cuspidal eigenform of parallel weight $\ell\ge2$, cannot
be a Hecke eigenform.
\end{proposition}

\begin{proof}
Put $w=(1+\sqrt{21})/2$, so $w^2-w-5=0$.
We have $h_F=1$, $h_F^+=2$, and $\zeta_F(-1)=1/3$.
Thus $\psi=\phi$, $\phi^2=1$, and the normalized hypothetical product
is $G=12E_2(\phi,\phi)h$. Write $a_m=c((m),h)$ and let $\eta$
be the central character of $h$.

The prime above $3$ is generated by $1+w$, whose norm is $-3$.
Hence $\eta(\mathfrak p_3)=(-1)^\ell$. Its norm is less than four,
so the small-norm coefficient identity gives
$c(\mathfrak p_3,G)=\phi(\mathfrak p_3)c(\mathfrak p_3,h)$.
Applying the two Hecke recurrences gives
\[
 c((3),G)-a_3=-8(-1)^\ell3^{\ell-1}.
\]
The only totally positive decompositions of $3$ are $1+2$ and $2+1$.
The Fourier product therefore gives the same difference as $12(a_2+5)$,
and consequently
\[
 a_2=-5-2(-1)^\ell3^{\ell-2}.
\]
Since $(2)$ is an inert prime of norm four, Ramanujan requires
$|a_2|\le2^\ell$. The displayed expression violates this for every
even $\ell\ge2$ and every odd $\ell\ge5$; the inequalities follow
from the base cases $\ell=2,4,5$ and induction in steps of two.
Only $\ell=3$ remains, and then $a_2=1$.

The coefficient comparison at $(4)$ gives
$3a_2+a_3=-1-5\cdot4^{\ell-2}$. Hence $a_3=-24$, while the Hecke
recurrence at $(2)$ gives $a_4=a_2^2-4^2=-15$.
There are exactly six ordered totally positive decompositions of $5$:
\[
 (1,4),(2,3),(3,2),(4,1),(2+w,3-w),(3-w,2+w).
\]
Indeed, for $x=a+bw$ with both $x$ and $5-x$ totally positive,
$|b|\sqrt{21}<5$, so $b\in\{-1,0,1\}$, and the stated list follows.
The last two elements are totally positive units. The Eisenstein
coefficients at $(2),(3),(4)$ are $5,13,21$, respectively, so
\[
 c((5),G)-a_5
 =12\bigl(a_4+5a_3+13a_2+21+2\bigr)=-1188.
\]

On the other hand, $(5)=\mathfrak p\mathfrak q$ splits. The generators
$w$ and $1-w$ have norm $-5$, so both prime ideals represent the
nontrivial narrow class. In this class there is no ideal of norm one
or two: the unit ideal is narrow-principal, and $(2)$ is inert.
If an exponent corresponding to either prime ideal decomposed in its
narrow component, the two summand ideals would therefore each have norm
at least three. The norm inequality would give
$\sqrt5\ge\sqrt3+\sqrt3$, which is impossible.
Thus both exponents are indecomposable and
\[
 c(\mathfrak p,G)=\phi(\mathfrak p)c(\mathfrak p,h),\qquad
 c(\mathfrak q,G)=\phi(\mathfrak q)c(\mathfrak q,h).
\]
Multiplicativity and $\phi((5))=1$ force $c((5),G)=a_5$,
contradicting the previous calculation.
\end{proof}
